%*******************************************************
% Mục 5 đến 10
%*******************************************************
\chapter{Lý do chọn cơ sở nghiên cứu}\label{ch:5}

Thí sinh chọn Trường Quốc tế, Đại học Quốc gia Hà Nội vì ba lý do. Thứ nhất, theo Thông báo tuyển sinh năm 2026, chương trình Tiến sĩ Kinh tế và Quản lý có tính liên ngành, tập trung vào kinh tế hành vi và quản lý trong lĩnh vực kinh tế; đề tài thuộc tài chính định lượng và giao thoa với quản lý ở phần định hướng chuyên môn hóa thị trường vốn Hà Nội.

Thứ hai, chương trình giảng dạy bằng tiếng Anh hoặc tiếng Việt. Học bằng tiếng Anh thuận lợi cho công bố quốc tế, và thí sinh có chứng chỉ tiếng Anh Bậc 4 theo Khung năng lực ngoại ngữ 6 bậc dùng cho Việt Nam.

Thứ ba, thí sinh tốt nghiệp ngành Kinh tế, chuyên ngành Đầu tư; theo đối chiếu của thí sinh, ngành này thuộc Nhóm 1 (Kinh tế học) trong Phụ lục 1 của Thông báo, nên thí sinh cần học bổ sung 9 tín chỉ. Nhà trường sẽ xác định điều kiện cuối cùng.

Ngoài ba lý do trên, đề tài gắn trực tiếp với Hà Nội: câu hỏi chính sách là bài toán số 19 của thành phố, sự kiện được nghiên cứu là việc sàn cổ phiếu niêm yết của Hà Nội chuyển vào HOSE, và người dùng dự kiến của kết quả là các cơ quan của thành phố và HNX. Một cơ sở đào tạo đặt tại Hà Nội thuận lợi cho việc làm việc với các đơn vị này trong suốt bốn năm.

\chapter{Mục tiêu và mong muốn đạt được khi đăng ký đi học nghiên cứu sinh}\label{ch:6}

Mục tiêu đầu tiên của thí sinh là đạt năng lực nghiên cứu độc lập ở trình độ tiến sĩ trong tài chính định lượng: thiết kế một nghiên cứu thực nghiệm có đăng ký trước, ước lượng tác động của một sự kiện duy nhất với suy luận phù hợp, và báo cáo trung thực cả những kết quả chưa đủ cơ sở kết luận. Thí sinh muốn rèn các kỹ năng này trong một môi trường có học phần phương pháp, hội thảo định kỳ và cơ chế phản biện nội bộ.

Mục tiêu thứ hai là tạo ra bằng chứng có ích cho hai nhóm người dùng: nhà nghiên cứu về trung tâm tài chính và cấu trúc thị trường; và Ủy ban nhân dân thành phố Hà Nội, cơ quan quản lý thị trường chứng khoán, HNX và các tổ chức đầu tư. Thí sinh mong kết quả của luận án là một đóng góp nhỏ, có kiểm chứng, cho câu hỏi số 1 và số 3 của bài toán số 19 (mục~\ref{sec:lydo}), đúng vào thời điểm thị trường vốn Hà Nội thay đổi cấu trúc.

Ba kỹ năng ở mục tiêu đầu tiên tương ứng với ba phần khó nhất của đề tài: đăng ký trước (mục~\ref{sec:mohinh}), suy luận khi chỉ có một sự kiện (mục~\ref{sec:kiemdinh}) và cách báo cáo kết quả chưa đủ cơ sở kết luận (mục~\ref{sec:tincay}). Đề tài vì vậy vừa là nội dung nghiên cứu, vừa là môi trường để rèn đúng những kỹ năng thí sinh muốn có.

\chapter{Những dự định và kế hoạch để đạt được những mục tiêu mong muốn}\label{ch:7}

\section{Học phần đào tạo trình độ tiến sĩ}

Trong hai quý đầu năm thứ nhất, thí sinh hoàn thành 9 tín chỉ học bổ sung (học phần Hành vi tổ chức và Lãnh đạo cùng hai học phần tự chọn) và các học phần theo chương trình, ưu tiên học phần tự chọn về kinh tế lượng chuỗi thời gian và phương pháp nghiên cứu định lượng.

Học phần kinh tế lượng chuỗi thời gian được ưu tiên vì phần suy luận chính của \HH{2} và \HH{3} dựa trên các mô hình chuỗi thời gian: bootstrap khối dừng, mô hình tự hồi quy và điểm gãy cấu trúc.

\section{Chuyên đề tiến sĩ}

Thí sinh dự kiến ba chuyên đề bám theo cấu trúc luận án: (i) tổng quan lý thuyết về trung tâm tài chính, cấu trúc sàn và hợp nhất sàn; (ii) chất lượng thị trường vốn của HNX, UPCoM và HOSE trước hợp nhất; (iii) tác động của hợp nhất sàn lên thanh khoản và khả năng đầu tư. Mỗi chuyên đề đứng độc lập được và làm nền cho một chương của luận án.

\section{Phần nghiên cứu khoa học}\label{sec:nckh}

Thí sinh dự kiến công bố ba bài báo trên tạp chí khoa học có phản biện và trình bày một đến hai báo cáo tại hội thảo khoa học. Bài báo 1 trình bày Chỉ số chất lượng thị trường vốn và so sánh HNX với HOSE (\HH{2}); bài báo 2 trình bày tác động của việc chuyển sàn lên thanh khoản và khả năng đầu tư (\HH{3}, \HH{4}); bài báo 3 trình bày bằng chứng quốc tế (\HH{1}) và so sánh vai trò của Hà Nội với các trung tâm trong khu vực. Mục tiêu là ít nhất một bài, do thí sinh là tác giả chính, thuộc cơ sở dữ liệu Web of Science hạng Q1 hoặc Q2.

\section{Phần thực hiện luận án}

Thời gian đào tạo đối với người có bằng đại học là bốn năm, theo hình thức tập trung toàn thời gian; thời gian này không vượt quá 48 tháng, là thời gian hỗ trợ tối đa theo Điều 13 Nghị quyết số 57/2026/NQ-HĐND (Hội đồng nhân dân thành phố Hà Nội, 2026). Bảng~\ref{tab:kehoach} nêu kế hoạch và sản phẩm dự kiến của từng năm.

\begin{table}[H]
\myfloatalign
\caption{Kế hoạch thực hiện luận án theo bốn năm và sản phẩm dự kiến}\label{tab:kehoach}
{\bangnho
\begin{tabularx}{\textwidth}{@{}l L >{\raggedright\arraybackslash}p{0.3\textwidth}@{}}
\toprule
Năm & Nội dung & Sản phẩm\\
\midrule
Năm 1 & Trước khi nhập học: đăng ký trước kế hoạch phân tích (trước ngày 28/12/2026). Quý 1: xác nhận nguồn và giấy phép dữ liệu HNX, UPCoM, HOSE, dữ liệu công bố và báo cáo GFCI. Quý 1--2: học bổ sung 9 tín chỉ và các học phần theo chương trình. Quý 2--3: làm sạch dữ liệu 2009--2026, tính các chỉ báo thành phần, tính hiệu ứng nhỏ nhất phát hiện được. Quý 4: xây dựng Chỉ số chất lượng thị trường vốn, kiểm định \HH{2} & Tệp đăng ký trước; tổng quan nghiên cứu; chuyên đề 1; Chỉ số chất lượng thị trường vốn bản đầu\\
\addlinespace
Năm 2 & Quý 1--2: viết bài báo 1 (\HH{2}); số hóa điểm GFCI và kiểm định \HH{1}. Quý 3--4: ước lượng giữa kỳ cho \HH{3} với 12 tháng dữ liệu sau sự kiện, báo cáo là kết quả giữa kỳ theo tệp đăng ký & Bài báo 1; kết quả giữa kỳ \HH{3}; báo cáo hội thảo; chuyên đề 2\\
\addlinespace
Năm 3 & Quý 1--2: hoàn tất 24 tháng dữ liệu sau sự kiện; kiểm định \HH{3} và ước lượng \HH{4}; viết bài báo 2 (\HH{3}, \HH{4}). Quý 3--4: so sánh Hà Nội với Bắc Kinh, Thành phố Hồ Chí Minh và các trung tâm khu vực; viết bài báo 3 (\HH{1}) & Bài báo 2; chỉ số khả năng đầu tư; chuyên đề 3\\
\addlinespace
Năm 4 & Quý 1--2: đề xuất định hướng chuyên môn hóa và chính sách cho thị trường vốn Hà Nội; công bố mã nguồn hai chỉ số. Quý 3--4: hoàn thiện luận án, chỉnh sửa theo góp ý và bảo vệ & Bài báo 3; báo cáo định hướng chuyên môn hóa gửi Sở Khoa học và Công nghệ Hà Nội; luận án; hai chỉ số mã nguồn mở\\
\bottomrule
\end{tabularx}}
\nguon{Nguồn: Thí sinh đề xuất}
\end{table}

Kế hoạch được sắp xếp theo một ràng buộc không đổi được: dữ liệu sau sự kiện chỉ có dần theo thời gian. Mốc 24 tháng sau sự kiện là hết tháng 12/2028, nên kiểm định \HH{3} sớm nhất là năm thứ ba. Hai năm đầu vì vậy dành cho những phần không phụ thuộc vào dữ liệu sau sự kiện: chỉ số chất lượng thị trường (\HH{2}) và bằng chứng quốc tế (\HH{1}). Kết quả giữa kỳ của \HH{3} ở năm thứ hai được ghi rõ là giữa kỳ, để không bị nhầm với kết quả theo kế hoạch đăng ký.

\chapter{Kiến thức chuyên môn của người dự tuyển tiến sĩ}\label{ch:8}

\section{Kinh nghiệm trong hoạt động chuyên môn}

Kinh nghiệm chuyên môn của thí sinh chủ yếu là kinh nghiệm nghiên cứu. Thí sinh chủ nhiệm đề tài nghiên cứu khoa học sinh viên \enquote{Nâng cao hiệu quả danh mục đầu tư bằng mô hình Black-Litterman: Áp dụng cho ngành ngân hàng trên thị trường chứng khoán Việt Nam}, đạt Giải Nhất cấp Khoa và Giải Nhì cấp Học viện năm 2026; đề tài đang dự Giải thưởng Khoa học và Công nghệ dành cho sinh viên năm 2026 do Bộ Giáo dục và Đào tạo tổ chức. Thí sinh còn đạt Giải Nhì nghiên cứu khoa học sinh viên cấp Khoa năm 2025 và Giải Khuyến khích cuộc thi Data Explorer. Bảng~\ref{tab:congtrinh} liệt kê các công trình đã công bố hoặc đã được chấp nhận đăng.

\begin{table}[H]
\myfloatalign
\caption{Công trình đã công bố hoặc đã được chấp nhận đăng}\label{tab:congtrinh}
{\bangnho
\begin{tabularx}{\textwidth}{@{}l L >{\raggedright\arraybackslash}p{0.27\textwidth} >{\raggedright\arraybackslash}p{0.18\textwidth}@{}}
\toprule
TT & Công trình & Tạp chí; ISSN & Trạng thái\\
\midrule
1 & Nguyễn Thanh Bình, Nguyễn Văn Trung*, \enquote{Tiền gửi không kỳ hạn, hiệu quả hoạt động và ổn định tài chính ngân hàng: Bằng chứng từ mô hình ngưỡng tại Việt Nam} (Nguyễn \& Nguyễn, 2026b) & Tạp chí Kinh tế - Luật và Ngân hàng; ISSN 3030-4199; Tập 28, Số 9, tr. 52--64; DOI 10.59276/3030-4199/jelb.bav.895 & Đã công bố\\
\addlinespace
2 & Nguyễn Thanh Bình, Nguyễn Văn Trung, \enquote{Black-Litterman portfolio optimization using regime switching CAPM and ABC-MCMC: Empirical evidence from the Vietnamese stock market period 2019 - 2025} (Nguyễn \& Nguyễn, 2026a) & Tạp chí Nghiên cứu Chính sách và Phát triển; ISSN 3030-4091; Tập 3, Số 1, tr. 81--99; DOI 10.63640/3030-4091/jpd.apd.194 & Đã công bố\\
\addlinespace
3 & Nguyễn Thanh Bình, Nguyễn Văn Trung*, Nguyễn Anh Tuấn, \enquote{GRI adoption and corporate brownwashing: Board governance evidence from ASEAN-5} & International Journal of Management and Sustainability (Conscientia Beam); ISSN 2306-9856 (bản in), 2306-0662 (bản trực tuyến); Scopus Q3 & Đã được chấp nhận đăng năm 2026, chưa xuất bản tại ngày 30/9/2026\\
\addlinespace
4 & Nguyen Thanh Binh, Nguyen Van Trung*, Nguyen Bach Diep, Ha Hong Hanh, Le Hong Minh, \enquote{Portfolio optimization with the inverse Black--Litterman framework and clustering machine-learning models: Evidence from Vietnamese bank stocks} & Multidisciplinary Science Journal (Malque Publishing); ISSN 2675-1240; Scopus Q3 & Đã được chấp nhận đăng, công bố trực tuyến trong mục Accepted Articles ngày 09/10/2026 (mã 19322; \url{https://malque.pub/ojs/index.php/msj/article/view/19322})\\
\bottomrule
\end{tabularx}}
\nguon{Ghi chú: (*) Tác giả liên hệ. Tên bài báo bằng tiếng Anh giữ nguyên như khi công bố. Trạng thái tính đến ngày 30/9/2026; riêng bài số 4 cập nhật đến ngày 09/10/2026.}
\end{table}

Ngoài các công trình trên, thí sinh có tám bản thảo đang ở các bước phản biện hoặc xét duyệt (Bảng~\ref{tab:banthao}), trong đó hai bản thảo liên quan trực tiếp đến đề tài.

\begin{table}[H]
\myfloatalign
\caption{Bản thảo đã gửi tạp chí, chưa công bố tại ngày 30/9/2026}\label{tab:banthao}
{\footnotesize\setlength{\tabcolsep}{3pt}
\begin{tabularx}{\textwidth}{@{}l L c >{\raggedright\arraybackslash}p{0.2\textwidth} >{\raggedright\arraybackslash}p{0.17\textwidth} >{\raggedright\arraybackslash}p{0.12\textwidth}@{}}
\toprule
TT & Tên bài báo & Số tác giả & Tạp chí; ISSN & Trạng thái (mã bài) & Liên quan đến đề tài\\
\midrule
1 & Tác động của cấu trúc sở hữu lên ổn định tài chính của các ngân hàng thương mại Việt Nam & 2 & Tạp chí Kinh tế - Luật và Ngân hàng; ISSN 3030-4199 & Qua 2 vòng phản biện, đang chỉnh sửa (mã 3129) & Không\\
2 & Limits to arbitrage in the Vietnamese physical gold market: The failure of short-term momentum & 3 & Journal of Economic and Banking Studies; ISSN 2734-9853 & Đang sửa vòng 1 (mã JEBS.1563) & Không\\
3 & Corporate Brownwashing and Firm Valuation in ASEAN-5 Disclosure Regimes & 2 & Cogent Business \& Management; ISSN 2331-1975; Scopus Q2 & Đang phản biện vòng 1 (mã 260865765) & Không\\
4 & Structural breaks in Vietnamese stock market liquidity: Evidence from the Amihud illiquidity measure and Bai-Perron regime-shift detection & 6 & Asian Journal of Economics and Banking; ISSN 2615-9821 & Chờ phân công phản biện (mã AJEB-08-2026-0208) & Có (chỉ báo thanh khoản, điểm gãy)\\
5 & Nested equity index correlations overstate true co-movement: Evidence from Vietnam & 4 & Asia-Pacific Financial Markets; ISSN 1387-2834; Scopus Q2 & Ban biên tập đang xem xét (mã FEJM-D-26-00726) & Không\\
6 & Who gains from a market upgrade? Stock liquidity and prices around Vietnam's FTSE Russell reclassification & 4 & Finance Research Open; ISSN 3050-7006 & Ban biên tập đang xem xét (mã FINR-D-26-00823) & Có (thanh khoản thị trường)\\
7 & Rated at the peak? Firm valuation around the first LSEG ESG score in five Southeast Asian markets & 4 & Sustainable Finance Review; ISSN 2977-7542 & Chờ xử lý tại tòa soạn (mã SFR-09-2026-0113) & Không\\
8 & Sales-based real earnings management and the cost of debt: Evidence from five ASEAN markets & 5 & Accounting Open; ISSN 3050-5461 & Ban biên tập đang xem xét (mã ACCOP-D-26-00232) & Không\\
\bottomrule
\end{tabularx}}
\nguon{Ghi chú: Thí sinh là tác giả liên hệ của cả tám bản thảo.}
\end{table}

\section{Kiến thức, sự hiểu biết và sự chuẩn bị của thí sinh về vấn đề nghiên cứu}

Thí sinh tốt nghiệp Cử nhân Kinh tế, chuyên ngành Đầu tư, hệ chính quy khóa 2022--2026 tại Học viện Chính sách và Phát triển, xếp loại Giỏi (điểm trung bình 3,41/4, tương đương 8,06/10). Các học phần gần nhất với đề tài gồm Quản lý danh mục đầu tư (9,9/10), Xác suất và thống kê (10/10), Kinh tế lượng (8,1/10) và Phân tích báo cáo tài chính (8,6/10).

Thí sinh đang theo học ngành Kế toán hệ vừa làm vừa học tại Đại học Kinh tế Quốc dân (năm thứ ba), ngành Công nghệ thông tin hệ đào tạo từ xa tại Học viện Công nghệ Bưu chính Viễn thông (năm thứ ba), và chương trình thạc sĩ trực tuyến tại WorldQuant University từ tháng 7/2026.

Nếu trúng tuyển, thí sinh cam kết theo học chương trình đào tạo tiến sĩ theo hình thức tập trung toàn thời gian, đáp ứng điều kiện tại khoản 3 Điều 11 Nghị quyết số 57/2026/NQ-HĐND (Hội đồng nhân dân thành phố Hà Nội, 2026), và trước khi nhập học sẽ hoàn thành, bảo lưu hoặc dừng các chương trình đang theo học nêu trên để dành toàn bộ thời gian cho học tập và nghiên cứu.

Về ngoại ngữ, thí sinh có chứng chỉ tiếng Anh Bậc 4 (VSTEP, Vietnamese Standardized Test of English Proficiency), điểm tổng 6,5, do Đại học Phenikaa cấp ngày 26/01/2026.

Thí sinh chuẩn bị cho đề tài qua các công việc nghiên cứu đã làm. Bản thảo số 6 ở Bảng~\ref{tab:banthao}, về thanh khoản cổ phiếu quanh đợt nâng hạng thị trường của FTSE Russell, cho thí sinh kinh nghiệm với nghiên cứu sự kiện trên thanh khoản, công cụ trực tiếp của \HH{3}; bản thảo số 4, về điểm gãy cấu trúc của thanh khoản, cho thí sinh kinh nghiệm với độ phi thanh khoản Amihud và phương pháp Bai-Perron. Bài báo Nguyễn và Nguyễn (2026a) là nền của mô hình \MM{5}, và bài báo số 4 ở Bảng~\ref{tab:congtrinh} giúp thí sinh làm quen với việc so sánh nhiều mô hình danh mục trên cùng dữ liệu.

Mối liên hệ giữa công việc đã làm và đề tài có thể tóm lại như sau: điểm gãy cấu trúc (bản thảo số 4) dùng để chọn cửa sổ trước của \HH{3}; nghiên cứu sự kiện trên thanh khoản (bản thảo số 6) là khung của phương trình \eqref{eq:es}; Black-Litterman có chuyển chế độ (Nguyễn \& Nguyễn, 2026a) là mô hình \MM{5}; so sánh nhiều mô hình danh mục (bài số 4 ở Bảng~\ref{tab:congtrinh}) là cách xây dựng chỉ số khả năng đầu tư.

\chapter{Dự kiến việc làm và các hướng nghiên cứu sau khi tốt nghiệp}\label{ch:9}

Sau khi tốt nghiệp, thí sinh dự kiến làm nghiên cứu và giảng dạy tài chính tại trường đại học hoặc viện nghiên cứu, hoặc làm việc ở bộ phận dữ liệu và phân tích thị trường của sở giao dịch hay cơ quan quản lý. Ở cả hai hướng, thí sinh có thể duy trì Chỉ số chất lượng thị trường vốn cùng người dùng thực tế.

Thí sinh cam kết thực hiện nghĩa vụ tại điểm e khoản 2 Điều 15 Nghị quyết số 57/2026/NQ-HĐND: sau khi hoàn thành chương trình đào tạo, tham gia mạng lưới chuyên gia khoa học, công nghệ, đổi mới sáng tạo và chuyển đổi số của Thành phố, và tham gia tư vấn, phản biện, cung cấp luận cứ khoa học cho bài toán số 19 và các nhiệm vụ khác phù hợp với chuyên môn khi được cơ quan có thẩm quyền mời hoặc giao nhiệm vụ (Hội đồng nhân dân thành phố Hà Nội, 2026).

Về nghiên cứu, thí sinh dự kiến mở rộng Chỉ số chất lượng thị trường vốn sang thị trường trái phiếu chính phủ và phái sinh do HNX vận hành, đánh giá tác động dài hạn của đợt hợp nhất, và so sánh vai trò của Hà Nội với các thủ đô khác trong khu vực.

\chapter{Đề xuất người hướng dẫn}\label{ch:10}

Thí sinh đề xuất hai người hướng dẫn. Người hướng dẫn chính thuộc định hướng nghiên cứu \enquote{chính sách tài khóa, tài chính tiền tệ, phát triển thị trường tài chính, tái cấu trúc hệ thống ngân hàng, quản trị tài chính, quản trị ngân hàng} trong danh mục của Nhà trường năm 2026; người hướng dẫn thứ hai là nhà khoa học ngoài Trường.

Người hướng dẫn chính (dự kiến): GS.TS. Nguyễn Đức Khương. Hướng nghiên cứu của người hướng dẫn gồm thị trường tài chính mới nổi, mô hình biến động và quản trị rủi ro trên thị trường vốn quốc tế, trùng với ba cấu phần của đề tài: cấu trúc và chất lượng thị trường vốn mới nổi, nghiên cứu sự kiện trên thị trường và quyết định danh mục.

Người hướng dẫn thứ hai (dự kiến): PGS.TS. Hà Hồng Hạnh, Trường Đại học Kinh tế Quốc dân. Người hướng dẫn hỗ trợ phần hàm ý chính sách và quản trị của đề tài trong bối cảnh khung quy định Việt Nam.
